\documentclass[11pt, a4paper]{article}

\usepackage{amsmath, amssymb, amsfonts}
\usepackage{graphicx}
\usepackage{geometry}
\usepackage{hyperref}
\usepackage{tikz}
\usepackage{booktabs}
\usepackage{xcolor}
\usepackage{bm}
\usepackage{float}
\usetikzlibrary{arrows.meta, positioning, fit, calc, shapes.geometric}

\geometry{margin=2.5cm}

\title{%
  \textbf{LQR with Optimal Feedback Control (OFC)} \\[0.5em]
  \large Mode: \texttt{lqr-applied-to-both-cart-manip} \\[0.2em]
  \normalsize 2D Welded Cart-Pendulum System with Muscle Dynamics and
  Joint-Space Manipulator Control
}
\author{Research Note}
\date{\today}

\begin{document}
\maketitle
\tableofcontents
\newpage

% ============================================================================
\section{Overview}
% ============================================================================

This note describes the formulation and implementation of the
\texttt{lqr-applied-to-both-cart-manip} simulation mode.
The key architectural difference from the
\texttt{lqr-applied-to-cart-manip-following-cart} mode is that the
\textbf{cart body is rigidly welded to the manipulator end-effector (EE)}.
There are therefore no prismatic actuators on the cart: the cart position is
\emph{entirely determined} by the manipulator's joint angles via forward
kinematics.

\begin{enumerate}
  \item \textbf{Manipulator (2~DOF):} A 2-link revolute arm whose EE is
        driven by an OFC chain (LQR $\to$ Muscle $\to$ ZFT $\to$ IK $\to$
        Computed-Torque) to reach the target $(x^*, y^*)$.
  \item \textbf{Cart-Pendulum (2~DOF):} A 3D gimbal pendulum welded to the
        manipulator EE. It receives no direct actuation — the pendulum is
        stabilised indirectly through the motion of the EE (and hence the
        cart body).
\end{enumerate}

The two sub-systems share a single Drake \texttt{MultibodyPlant} with
\emph{separate model instances} but a rigid weld constraint connecting the
cart body to the EE frame.

\medskip
\noindent\textbf{Key contrast with the other mode:}

\begin{center}
\renewcommand{\arraystretch}{1.3}
\begin{tabular}{lll}
\toprule
\textbf{Aspect} & \textbf{cart-manip-following-cart} & \textbf{both-cart-manip} \\
\midrule
Cart actuation    & Prismatic (impedance force)  & None (welded to EE) \\
Manipulator role  & Passive tracker via IK + CTC & Active driver of cart \\
Impedance block   & \checkmark\ present           & $\times$\ removed \\
Cart $[x,y]$      & Plant state (prismatic DOF)  & FK output: $\mathrm{FK}(q_1,q_2)$ \\
LQR state space   & 14D (cart + pend + muscle + ZFT) & 14D (arm + pend + muscle + ZFT) \\
\bottomrule
\end{tabular}
\end{center}

% ============================================================================
\section{State and Coordinate Definitions}
% ============================================================================

\subsection{Plant State (8D)}

Because the cart is welded to the EE, the plant no longer contains prismatic
DOF for the cart.  Instead the plant state is:

\[
\bm{s}_{\mathrm{plant}} =
\underbrace{[q_1,\; q_2]^{\top}}_{\text{arm joints}}
\oplus
\underbrace{[\alpha,\; \beta]^{\top}}_{\text{pend.\ angles}}
\oplus
\underbrace{[\dot{q}_1,\; \dot{q}_2]^{\top}}_{\text{arm vel.}}
\oplus
\underbrace{[\dot{\alpha},\; \dot{\beta}]^{\top}}_{\text{pend.\ ang.\ vel.}}
\in \mathbb{R}^{8}
\]

where $q_1$ is the base-to-link1 joint angle and $q_2$ is the link1-to-link2
joint angle.  Note that \textbf{Drake returns joints in URDF parse order},
which for this URDF is $[q_2, q_1]$; index conversion is handled internally.

\subsection{EE / Cart Position (computed)}

The cart body position in the world frame is \emph{not} a state variable; it
is computed from the arm state via forward kinematics:

\[
\bm{p} = [x, y]^{\top} = f_{\mathrm{FK}}(q_1, q_2) \in \mathbb{R}^{2}
\]
\[
\dot{\bm{p}} = J(q_1, q_2)\,[\dot{q}_1, \dot{q}_2]^{\top}, \quad
J = \frac{\partial f_{\mathrm{FK}}}{\partial \bm{q}} \in \mathbb{R}^{2\times 2}
\]

A dedicated Drake \texttt{LeafSystem} (\texttt{EndEffectorKinematics2D})
evaluates these at every integration step and outputs $(\bm{p}, \dot{\bm{p}})$.

\subsection{Full OFC / Augmented State (14D)}

The LQR operates on the augmented state:

\[
\bm{x} =
\begin{bmatrix}
q_1,\; q_2,\; \alpha,\; \beta,\;
\dot{q}_1,\; \dot{q}_2,\; \dot{\alpha},\; \dot{\beta},\;
F_x,\; F_y,\;
x_{\mathrm{ref}},\; y_{\mathrm{ref}},\;
\dot{x}_{\mathrm{ref}},\; \dot{y}_{\mathrm{ref}}
\end{bmatrix}^{\top} \in \mathbb{R}^{14}
\]

where the first 8 entries come from the plant, the next 2 from the muscle
dynamics, and the last 4 from the ZFT reference mass.

\subsection{Manipulator State (4D)}

\[
\bm{s}_{\text{manip}} =
[q_1,\; q_2,\; \dot{q}_1,\; \dot{q}_2]^{\top} \in \mathbb{R}^{4}
\]

% ============================================================================
\section{Control Architecture}
% ============================================================================

\subsection{Signal Chain}

All actuation flows through the manipulator joint torques:

\begin{center}
\textbf{LQR} $\xrightarrow{\bm{u}}$
\textbf{Muscle} $\xrightarrow{\bm{F}}$
\textbf{ZFT} $\xrightarrow{\bm{s}_{\mathrm{ref}}}$
\textbf{IK} $\xrightarrow{\bm{q}_d, \dot{\bm{q}}_d, \ddot{\bm{q}}_d}$
\textbf{CTC} $\xrightarrow{\bm{\tau}}$
\textbf{Manipulator Plant} $\xrightarrow{\bm{p} = f_{\mathrm{FK}}}$ (cart moves)
\end{center}

\noindent The pendulum angles $\alpha, \beta$ are shaped
\emph{indirectly} by the cart (EE) acceleration that results from the joint
torques — there is no direct torque input to the pendulum.

\subsection{Block Descriptions}

\subsubsection*{Block 1 — Muscle Dynamics}

\[
\boxed{
\dot{F}_i = \frac{-F_i + u_i}{\tau_m}, \quad i \in \{x, y\}
}
\]

\begin{itemize}
  \item \textbf{Input:} $\bm{u} \in \mathbb{R}^2$ (LQR neural command)
  \item \textbf{Output:} $\bm{F} = [F_x, F_y]^{\top} \in \mathbb{R}^2$ (muscle force)
  \item \textbf{Time constant:} $\tau_m = 0.03\;\mathrm{s}$
\end{itemize}

\subsubsection*{Block 2 — ZFT Reference Mass}

The muscle force drives a virtual reference mass that provides the smooth EE
target trajectory.  Its spring-damper coupling is to the \emph{EE position}
$\bm{p} = f_{\mathrm{FK}}(\bm{q})$, not to the plant cart state (there is no
prismatic cart state).

\[
\boxed{
M_{\mathrm{ref}}\,\ddot{\bm{p}}_{\mathrm{ref}}
= \bm{F}
+ K\,(\bm{p} - \bm{p}_{\mathrm{ref}})
+ D\,(\dot{\bm{p}} - \dot{\bm{p}}_{\mathrm{ref}})
}
\]

\begin{itemize}
  \item \textbf{Inputs:} $\bm{F}$ (muscle); $\bm{p}, \dot{\bm{p}}$ (EE kinematics via FK)
  \item \textbf{Output:} $\bm{s}_{\mathrm{ref}} = [x_{\mathrm{ref}}, y_{\mathrm{ref}},
        \dot{x}_{\mathrm{ref}}, \dot{y}_{\mathrm{ref}}]^{\top} \in \mathbb{R}^4$;
        also $\ddot{\bm{p}}_{\mathrm{ref}} \in \mathbb{R}^2$ for feed-forward
  \item \textbf{Parameters:} $M_{\mathrm{ref}} = 1.0\;\mathrm{kg}$,
        $K = 50\;\mathrm{N/m}$, $D = 10\;\mathrm{N\cdot s/m}$
\end{itemize}

\subsubsection*{Block 3 — IK Solver (ZFT Joint Reference)}

Rather than a one-shot IK per step, a \emph{differential IK} integrates joint
velocity commands resulting from the ZFT output:

\[
\boxed{
\dot{\bm{q}}_d = J^{\dagger}(\bm{q})\,
\bigl[\,K_p(\bm{p}_{\mathrm{ref}} - \bm{p}) + \dot{\bm{p}}_{\mathrm{ref}}\,\bigr]
+ J^{\dagger}(\bm{q})\,\ddot{\bm{p}}_{\mathrm{ref}}\,\Delta t
}
\]

where $J^{\dagger}$ is the Moore--Penrose pseudoinverse of the $2\times 2$
Jacobian.  The joint position reference $\bm{q}_d$ is obtained by integration.
This approach avoids the overhead of a full nonlinear IK solve at every step.

\begin{itemize}
  \item \textbf{Inputs:} $\bm{p}_{\mathrm{ref}}, \dot{\bm{p}}_{\mathrm{ref}},
        \ddot{\bm{p}}_{\mathrm{ref}}$ (ZFT outputs); plant state (for $J$ and warm-start)
  \item \textbf{Outputs:} $\bm{q}_d, \dot{\bm{q}}_d, \ddot{\bm{q}}_d \in \mathbb{R}^2$
\end{itemize}

\subsubsection*{Block 4 — Finite-Horizon LQR Controller}

The LQR minimises the quadratic cost over a finite horizon $T$:

\[
\boxed{
J = \int_0^T \left[\bm{e}^{\top} Q\, \bm{e}
                 + \bm{u}^{\top} R\, \bm{u}\right] dt
  + \bm{e}(T)^{\top} Q_N\, \bm{e}(T)
}
\]

where $\bm{e} = \bm{x} - \bm{x}^{*}$ is the 14D state error from the goal:

\[
\bm{x}^{*} =
[q_1^*, q_2^*, 0, 0, 0, 0, 0, 0, 0, 0, x^*, y^*, 0, 0]^{\top}
\]

with $(q_1^*, q_2^*)$ solved analytically from the target $(x^*, y^*)$ via
2-link IK.  The backward Riccati ODE gives the time-varying gain:

\[
\bm{u}(t) = -R^{-1}B^{\top}P(t)\,\bm{e}(t)
\]

\paragraph{Cost matrices:}
\begin{align*}
Q &= \mathrm{diag}(100, 100, 500, 500, 50, 50, 100, 100, 0.1, 0.1, 1, 1, 1, 1) \\
Q_N &= 2\,Q \\
R &= \mathrm{diag}(10, 10)
\end{align*}

The larger $R$ (compared to the other mode) penalises neural effort more
heavily, giving smoother joint motion.

\subsubsection*{Block 5 — Joint-Space Computed-Torque Controller}

\[
\boxed{
\bm{\tau} = M(\bm{q})\,\ddot{\bm{q}}_d
          + C(\bm{q}, \dot{\bm{q}})\,\dot{\bm{q}}
          + G(\bm{q})
          + K_p (\bm{q}_d - \bm{q})
          + K_d (\dot{\bm{q}}_d - \dot{\bm{q}})
}
\]

\begin{itemize}
  \item \textbf{Inputs:} $[\bm{q}_d, \dot{\bm{q}}_d, \ddot{\bm{q}}_d]$ (IK);
        $[\bm{q}, \dot{\bm{q}}]$ (plant)
  \item \textbf{Output:} $\bm{\tau} \in \mathbb{R}^2$ (joint torques, capped at
        $\pm 100\;\mathrm{N\cdot m}$)
  \item \textbf{Gains:} $K_p = 200\;\mathrm{N\cdot m/rad}$,
        $K_d = 60\;\mathrm{N\cdot m\cdot s/rad}$
\end{itemize}

The torques $\bm{\tau}$ are the \textbf{only} actuation input to the plant.

% ============================================================================
\section{Block Diagram}
% ============================================================================

\begin{figure}[H]
\centering
\resizebox{\textwidth}{!}{%
\begin{tikzpicture}[
  box/.style   = {draw, rounded corners=3pt, minimum width=2.8cm, minimum height=1.1cm,
                  text centered, font=\small\bfseries, fill=white},
  greenbox/.style = {box, fill=green!15},
  bluebox/.style  = {box, fill=blue!10},
  graybox/.style  = {box, fill=gray!12},
  arr/.style   = {-{Latex[length=2mm]}, thick},
  darr/.style  = {-{Latex[length=2mm]}, thick, dashed, gray},
  node distance = 1.6cm and 2.0cm,
]

% ---------------------------------------------------------------
% TOP ROW: State label + equations
% ---------------------------------------------------------------
\node[graybox, minimum width=6.0cm, minimum height=1.4cm, align=left, font=\small]
  (state_box) at (0, 3.8)
  {$\bm{x} = [q_1, q_2, \alpha, \beta, \dot{q}_1, \dot{q}_2, \dot{\alpha}, \dot{\beta},$\\
   $\phantom{\bm{x}=}F_x, F_y, x_{\mathrm{ref}}, y_{\mathrm{ref}}, \dot{x}_{\mathrm{ref}}, \dot{y}_{\mathrm{ref}}]^\top \in \mathbb{R}^{14}$};

\node[graybox, minimum width=7.5cm, minimum height=1.4cm, align=left, font=\small]
  (cost_box) at (9.5, 3.8)
  {$J = \int_0^T\!\bigl[\bm{e}^\top Q \bm{e} + \bm{u}^\top R \bm{u}\bigr]dt + \bm{e}(T)^\top Q_N \bm{e}(T)$\\
   $\bm{e} = \bm{x} - \bm{x}^*$, \quad $\bm{u}(t) = -R^{-1}B^\top P(t)\bm{e}(t)$};

% ---------------------------------------------------------------
% MAIN SIGNAL PATH (middle row)
% ---------------------------------------------------------------
\node[bluebox]  (lqr)  at (0, 1.5)   {LQR\\Controller};
\node[box]      (musc) at (3.0, 1.5) {Muscle\\Dynamics};
\node[box]      (zft)  at (6.2, 1.5) {ZFT Reference\\Mass};
\node[box]      (ik)   at (9.4, 1.5) {Differential\\IK Solver};
\node[box]      (ctc)  at (12.6, 1.5){Computed-\\Torque Ctrl};
\node[greenbox] (plant) at (15.8, 1.5){Manipulator\\Plant\\(+ Welded Cart)};

% Arrows: main path
\draw[arr] (lqr)  -- node[above,font=\tiny] {$\bm{u}$} (musc);
\draw[arr] (musc) -- node[above,font=\tiny] {$\bm{F}$} (zft);
\draw[arr] (zft)  -- node[above,font=\tiny] {$\bm{s}_{\mathrm{ref}}$\\$\ddot{\bm{p}}_{\mathrm{ref}}$} (ik);
\draw[arr] (ik)   -- node[above,font=\tiny] {$\bm{q}_d,\dot{\bm{q}}_d,\ddot{\bm{q}}_d$} (ctc);
\draw[arr] (ctc)  -- node[above,font=\tiny] {$\bm{\tau}$} (plant);

% ---------------------------------------------------------------
% EE Kinematics block (below ZFT)
% ---------------------------------------------------------------
\node[graybox, minimum width=2.8cm, minimum height=0.9cm, font=\small]
  (eekin) at (6.2, -0.6)
  {EE Kinematics\\$\bm{p}=f_{\mathrm{FK}}(\bm{q})$};

% Plant feeds EE kin
\draw[darr] (plant.south) -- ++(0,-0.4) -| node[right, font=\tiny, pos=0.75] {$\bm{q}, \dot{\bm{q}}$} (eekin.east);

% EE kin feeds ZFT (EE position as cart position)
\draw[arr] (eekin.north) -- node[right,font=\tiny] {$\bm{p},\dot{\bm{p}}$} (zft.south);

% ---------------------------------------------------------------
% Equations below blocks
% ---------------------------------------------------------------
\node[font=\scriptsize, below=0.15cm of musc, align=center]
  {$\dot{F}_i = \tfrac{-F_i+u_i}{\tau_m}$};

\node[font=\scriptsize, below=0.15cm of zft, yshift=-0.9cm, align=center]
  {$M_{\mathrm{ref}}\ddot{\bm{p}}_{\mathrm{ref}} = \bm{F} + K(\bm{p}{-}\bm{p}_{\mathrm{ref}}) + D(\dot{\bm{p}}{-}\dot{\bm{p}}_{\mathrm{ref}})$};

\node[font=\scriptsize, below=0.15cm of ctc, align=center]
  {$\bm{\tau} = M\ddot{\bm{q}}_d + C\dot{\bm{q}} + G + K_p(\bm{q}_d{-}\bm{q}) + K_d(\dot{\bm{q}}_d{-}\dot{\bm{q}})$};

% ---------------------------------------------------------------
% 14D state feedback path
% ---------------------------------------------------------------
% Plant -> state mux (shown as curved path going below)
\draw[darr] (plant.south) ++(0,-0.4) -- ++(0,-1.4)
    -- ++(-17.8, 0)
    -- (lqr.south |- {coordinate (0,-2.4)})
    -- (lqr.south)
    node[below, font=\tiny, pos=0.35] {$\bm{x}\ (14\mathrm{D},\ \text{via ZOH})$};

% muscle -> aug state (implicit, shown as text annotation)
\node[font=\scriptsize, below=2.3cm of lqr, align=center, text width=2.5cm]
  {Full OFC State (14D)};

% ---------------------------------------------------------------
% IK also takes plant state for warm-start
% ---------------------------------------------------------------
\draw[darr] (plant.south) ++(0,-0.4) -- ++(0,-0.8)
    -- (ik.south |- {coordinate (0,-1.2)})
    -- (ik.south)
    node[right, font=\tiny, pos=0.6] {$\bm{q},\dot{\bm{q}}$};

% ---------------------------------------------------------------
% "No impedance" annotation
% ---------------------------------------------------------------
\node[font=\scriptsize\itshape, text=red!60!black, align=center]
  at (6.2, 0.45) {No impedance block:\\cart driven only by $\bm{\tau}$};

\end{tikzpicture}
}%
\caption{%
  Block diagram of the \texttt{lqr-applied-to-both-cart-manip} mode.
  The cart body is welded to the manipulator EE, so there are no prismatic
  actuators.  All actuation is through joint torques $\bm{\tau}$.
  Dashed arrows: feedback and warm-start signals.
  The EE kinematics block converts arm joint angles to task-space position
  $\bm{p} = f_{\mathrm{FK}}(\bm{q})$, which replaces the cart state in the
  ZFT ODE.}
\label{fig:block_both}
\end{figure}

% ============================================================================
\section{Signal Flow Summary}
% ============================================================================

\begin{table}[H]
\centering
\caption{Signal connections in the control loop (\texttt{lqr-applied-to-both-cart-manip}).}
\label{tab:signals}
\renewcommand{\arraystretch}{1.3}
\begin{tabular}{llll}
\toprule
\textbf{From} & \textbf{To} & \textbf{Signal} & \textbf{Dim.} \\
\midrule
LQR Controller        & Muscle Dynamics       & $\bm{u}$ (neural command)             & $\mathbb{R}^2$ \\
Muscle Dynamics       & ZFT Reference Mass    & $\bm{F}$ (muscle force)               & $\mathbb{R}^2$ \\
Muscle Dynamics       & State Mux (14D)       & $\bm{F}$ (for LQR state)             & $\mathbb{R}^2$ \\
EE Kinematics         & ZFT Reference Mass    & $\bm{p}, \dot{\bm{p}}$ (EE as cart)  & $\mathbb{R}^4$ \\
ZFT Reference Mass    & IK Solver             & $\bm{p}_{\mathrm{ref}}, \dot{\bm{p}}_{\mathrm{ref}}, \ddot{\bm{p}}_{\mathrm{ref}}$ & $\mathbb{R}^6$ \\
ZFT Reference Mass    & State Mux (14D)       & $\bm{s}_{\mathrm{ref}}$               & $\mathbb{R}^4$ \\
IK Solver             & Computed-Torque Ctrl  & $\bm{q}_d, \dot{\bm{q}}_d, \ddot{\bm{q}}_d$ & $\mathbb{R}^6$ \\
Manipulator Plant     & Computed-Torque Ctrl  & $[\bm{q}, \dot{\bm{q}}]$ (arm state) & $\mathbb{R}^4$ \\
Computed-Torque Ctrl  & Manipulator Plant     & $\bm{\tau}$ (joint torques)           & $\mathbb{R}^2$ \\
Manipulator Plant     & EE Kinematics         & $\bm{q}, \dot{\bm{q}}$ (plant state) & $\mathbb{R}^4$ \\
Manipulator Plant     & IK Solver             & plant state (warm-start)              & $\mathbb{R}^8$ \\
\midrule
Manipulator Plant     & State Mux             & $\bm{s}_{\mathrm{plant}}$ (8D)       & $\mathbb{R}^8$ \\
State Mux             & ZOH ($10\;\mathrm{ms}$) & $\bm{x}$ (14D)                     & $\mathbb{R}^{14}$ \\
ZOH                   & LQR Controller        & $\hat{\bm{x}}$ (held)                & $\mathbb{R}^{14}$ \\
\bottomrule
\end{tabular}
\end{table}

% ============================================================================
\section{Linearised System}
% ============================================================================

The LQR is designed on the 14D augmented linearisation obtained from the
welded plant.  The equilibrium $(\bm{x}^*, \bm{u}^* = \bm{0})$ is:

\[
\bm{x}^* = [q_1^*, q_2^*, 0, 0, 0, 0, 0, 0, 0, 0, x^*, y^*, 0, 0]^{\top}
\]

where $(q_1^*, q_2^*)$ is solved analytically from the 2-link inverse kinematics
for the target $(x^*, y^*)$.  Around this equilibrium:

\[
\dot{\bm{x}} = A\,\bm{x} + B\,\bm{u}
\]

The matrices $A \in \mathbb{R}^{14\times 14}$ and $B \in \mathbb{R}^{14\times 2}$
are obtained by \texttt{build\_linearized\_for\_complete\_system\_2d()}, which
uses Drake's automatic Jacobian (via \texttt{Linearize}) on the welded plant
augmented with muscle and ZFT dynamics.

The Riccati equation is solved backward in time from $t = T$ to $t = 0$ at
$\Delta t = 10\;\mathrm{ms}$ to obtain the gain schedule
$K(t) = R^{-1}B^{\top}P(t)$, which is interpolated online.

% ============================================================================
\section{Simulation Parameters}
% ============================================================================

\begin{table}[H]
\centering
\caption{Default simulation parameters.}
\renewcommand{\arraystretch}{1.3}
\begin{tabular}{lll}
\toprule
\textbf{Parameter} & \textbf{Value} & \textbf{Description} \\
\midrule
Arm link 1 length   & $\approx 2.20\;\mathrm{m}$ & From URDF geometry \\
Arm link 2 length   & $\approx 1.25\;\mathrm{m}$ & From URDF geometry \\
Pendulum mass       & $0.3\;\mathrm{kg}$         & \\
Pendulum length     & $0.25\;\mathrm{m}$         & \\
Arm damping         & $0.1\;\mathrm{N\cdot s/rad}$ & Both joints \\
Pendulum damping    & $0.1\;\mathrm{N\cdot s/rad}$ & Both gimbal axes \\
$\tau_m$            & $0.03\;\mathrm{s}$          & Muscle time constant \\
$M_{\mathrm{ref}}$  & $1.0\;\mathrm{kg}$          & ZFT reference mass \\
$K$                 & $50\;\mathrm{N/m}$          & ZFT impedance stiffness \\
$D$                 & $10\;\mathrm{N\cdot s/m}$   & ZFT impedance damping \\
$K_p$               & $200\;\mathrm{N\cdot m/rad}$        & CTC position gain \\
$K_d$               & $60\;\mathrm{N\cdot m\cdot s/rad}$  & CTC velocity gain \\
$\tau_{\max}$       & $100\;\mathrm{N\cdot m}$    & Joint torque limit \\
LQR horizon $T$     & $10.0\;\mathrm{s}$          & \\
$\Delta t$ (LQR)    & $10\;\mathrm{ms}$           & Riccati integration step \\
Simulation $\Delta t$ & $1\;\mathrm{ms}$          & Drake integrator step \\
\bottomrule
\end{tabular}
\end{table}

% ============================================================================
\section{Implementation Notes}
% ============================================================================

\begin{itemize}
  \item \textbf{Weld constraint}: The cart body is added to the plant inside
        \texttt{SystemBuilder.build()} when \texttt{weld\_cart\_to\_manip\_ee=True}.
        Drake enforces the constraint algebraically so the cart position is
        always equal to the EE pose with zero constraint force residual.

  \item \textbf{No impedance block}: Unlike the companion mode, the impedance
        force block (\texttt{ImpedanceForce2D}) is \textbf{not instantiated}.
        The sole actuation path is joint torques $\bm{\tau}$.

  \item \textbf{Differential IK}: The \texttt{ZFTJointReferenceIK} system uses
        the velocity-integration method (\texttt{ik\_method="differential"})
        with position feedback gain $K_p = 10\;\mathrm{s}^{-1}$ rather than
        solving a full nonlinear IK problem at every step.  This avoids
        kinematic singularities and is faster.

  \item \textbf{EE kinematics}: The \texttt{EndEffectorKinematics2D} system
        computes $(\bm{p}, \dot{\bm{p}})$ from the plant context at every
        step using Drake's \texttt{CalcPointsPositions} and the geometric
        Jacobian.

  \item \textbf{Drake joint ordering}: The URDF returns joints in parse order
        $[q_2, q_1]$.  All user-facing code uses $[q_1, q_2]$ order;
        conversions happen inside \texttt{CupManipulator}.

  \item \textbf{ZOH}: A Zero-Order Hold at $10\;\mathrm{ms}$ on the 14D
        augmented state breaks the algebraic loop between the LQR and the
        continuous plant/muscle/ZFT dynamics.
\end{itemize}

% ============================================================================
\newpage
\section{Equation Flow Summary}
% ============================================================================

This section traces the complete signal chain from sensor reading to actuator
output at every simulation step.

\subsection*{Step-by-step signal chain}

\bigskip
\noindent\textbf{(1) Read plant state.}
The plant outputs its 8D state:
\[
\bm{s}_{\mathrm{plant}}(t)
= [\,q_1,\;q_2,\;\alpha,\;\beta,\;\dot{q}_1,\;\dot{q}_2,\;\dot{\alpha},\;\dot{\beta}\,]^{\!\top}
\in\mathbb{R}^{8}
\]

\bigskip
\noindent\textbf{(2) EE kinematics — compute cart position via FK.}
The arm joint angles $\bm{q} = [q_1, q_2]^{\top}$ are passed to the EE
kinematics block:
\[
\bm{p}(t) = f_{\mathrm{FK}}(\bm{q}),
\qquad
\dot{\bm{p}}(t) = J(\bm{q})\,\dot{\bm{q}}
\]
This $(\bm{p}, \dot{\bm{p}})$ substitutes for the cart state in all downstream
blocks — the welded cart has no independent position state.

\bigskip
\noindent\textbf{(3) Muscle dynamics — low-pass filter on neural command.}
\[
\dot{\bm{F}}(t) = \frac{-\bm{F}(t) + \bm{u}(t)}{\tau_m}
\quad\Longrightarrow\quad
\bm{F}(t) \in\mathbb{R}^{2}
\]
$\bm{F}$ is fed to the ZFT reference mass and also stacked into the 14D state.

\bigskip
\noindent\textbf{(4) ZFT reference mass — smooth end-effector target.}
\[
\ddot{\bm{p}}_{\mathrm{ref}}(t)
= \frac{1}{M_{\mathrm{ref}}}\!\left[
    \bm{F}(t)
    + K\bigl(\bm{p}(t) - \bm{p}_{\mathrm{ref}}(t)\bigr)
    + D\bigl(\dot{\bm{p}}(t) - \dot{\bm{p}}_{\mathrm{ref}}(t)\bigr)
  \right]
\]
Integration gives:
\[
\bm{s}_{\mathrm{ref}}(t)
= [\,x_{\mathrm{ref}},\;y_{\mathrm{ref}},\;\dot{x}_{\mathrm{ref}},\;\dot{y}_{\mathrm{ref}}\,]^{\!\top}
\in\mathbb{R}^{4}
\]
All three outputs ($\bm{p}_{\mathrm{ref}},\dot{\bm{p}}_{\mathrm{ref}},
\ddot{\bm{p}}_{\mathrm{ref}}$) feed the IK solver.

\bigskip
\noindent\textbf{(5) Full 14D state assembly.}
\[
\bm{x}(t) =
\underbrace{[q_1, q_2, \alpha, \beta]}_{\text{joint/pend pos.}}
\oplus
\underbrace{[\dot{q}_1, \dot{q}_2, \dot{\alpha}, \dot{\beta}]}_{\text{joint/pend vel.}}
\oplus
\underbrace{[F_x, F_y]}_{\text{muscle}}
\oplus
\underbrace{[x_{\mathrm{ref}}, y_{\mathrm{ref}}, \dot{x}_{\mathrm{ref}}, \dot{y}_{\mathrm{ref}}]}_{\text{ZFT ref.}}
\]
A ZOH ($\Delta = 10\;\mathrm{ms}$) samples this to $\hat{\bm{x}}$.

\bigskip
\noindent\textbf{(6) Finite-horizon LQR — optimal neural command.}
\[
\bm{u}(t) = -\underbrace{R^{-1}B^{\!\top}P(t)}_{K(t)}\,\bigl(\hat{\bm{x}}(t) - \bm{x}^{*}\bigr)
\in\mathbb{R}^{2}
\]
This closes the outer loop back to Step 3 (muscle input) at the next step.

\bigskip
\noindent\textbf{(7) Differential IK — ZFT target to joint references.}
\[
\dot{\bm{q}}_d(t)
= J^{\dagger}(\bm{q})\bigl[K_p(\bm{p}_{\mathrm{ref}} - \bm{p})
+ \dot{\bm{p}}_{\mathrm{ref}}\bigr]
+ J^{\dagger}(\bm{q})\,\ddot{\bm{p}}_{\mathrm{ref}}\,\Delta t
\]
\[
\bm{q}_d(t) = \bm{q}_d(t{-}\Delta t) + \dot{\bm{q}}_d(t)\,\Delta t
\]

\bigskip
\noindent\textbf{(8) Computed-torque controller — joint torques.}
\[
\bm{\tau}(t)
= M(\bm{q})\ddot{\bm{q}}_d
+ C(\bm{q},\dot{\bm{q}})\dot{\bm{q}}
+ G(\bm{q})
+ K_p(\bm{q}_d - \bm{q})
+ K_d(\dot{\bm{q}}_d - \dot{\bm{q}})
\in\mathbb{R}^{2}
\]
$\bm{\tau}$ is the \textbf{sole actuation signal} applied to the welded plant.
The joint motion moves the EE (and the welded cart) toward the target, while
the pendulum is stabilised indirectly through the resulting EE acceleration.

\subsection*{Compact causal chain}

\[
\bm{s}_{\mathrm{plant}}
\;\xrightarrow{\text{FK}}\;
\bm{p}, \dot{\bm{p}}
\;\xrightarrow[\text{ZFT}]{\bm{F}}\;
\bm{p}_{\mathrm{ref}},\dot{\bm{p}}_{\mathrm{ref}},\ddot{\bm{p}}_{\mathrm{ref}}
\;\xrightarrow{\text{diff.\ IK}}\;
\bm{q}_d, \dot{\bm{q}}_d, \ddot{\bm{q}}_d
\;\xrightarrow{\text{CTC}}\;
\bm{\tau}
\;\longrightarrow\;
\text{plant}
\]
\[
\bm{s}_{\mathrm{plant}} \oplus \bm{F} \oplus \bm{s}_{\mathrm{ref}}
\;\xrightarrow{\text{mux}}\;
\bm{x}_{14}
\;\xrightarrow{\text{ZOH}}\;
\hat{\bm{x}}
\;\xrightarrow{\text{LQR}}\;
\bm{u}
\;\xrightarrow{\tau_m}\;
\bm{F}
\quad\text{(closes loop)}
\]

\noindent Each arrow represents a Drake system connection.  There is no
separate manipulator tracking loop — in this mode the manipulator \emph{is}
the cart actuator.

\end{document}
